\documentclass{subfiles}
\begin{document}
    At least in principle, hashing is very simple: 
        we use function \(h : K \to \set{0, 1, m - 1}\), where \(m = \abs{T}\),
        to map each of the keys in \(k\) to one of the slots in \(T\).

    It comes with no surprise that any two keys \(k, l \in K\) can be such that \(h(k) = h(l)\).
        If this happens, we say that a collision has occured. 
        Therefore we need a way to handle collisions. The literature provides several solutions,
        the simplest of which is \emph{chaining}. Under chaining each slot \(i\) points to a list \(L_{i}\),
        which in turns stores all the keys in \(K\) that maps to \(i\).
        In other words, for each \(i\)
        \[
            L_{i} = \set{k \in K \text{ s.t. } h(k) = i}.
        \]

    From what just discussed,
        it follows that we are interested in finding hash functions to minimise collisions.
        A class of such functions is discussed in the following section.
\end{document}
